\section{Balanced decisions: equivalence, conflict, and dynamics}
\label{app:decision-proofs}

This appendix supplies proofs and exact certificates for
Section~\ref{sec:sharing-dynamics}. Every partition has $K$ groups of
size $m=L/K$. All displayed layer indices are one-based. For row vectors
$x_i$ and a group $U$ of size $m$,
\begin{equation}
 \sum_{i\in U}\|x_i-\bar x_U\|^2
 =\frac1m\sum_{i<j,\ i,j\in U}\|x_i-x_j\|^2. \label{eq:pairwise-sse}
\end{equation}
This identity follows by expanding both sides; it also provides finite
integer certificates when all rows have a common rational denominator.

\subsection{Equivalence cases}

\begin{proposition}[Two-basis equivalence]
\label{prop:two-basis-decisions}
Let $K=2$, $L=2m$, and $b_1\ne b_2$. The globally optimal unlabeled
partitions for maximizing $S_A$ and minimizing $J_\Theta$ coincide,
including ties.
\end{proposition}
\begin{proof}
Write $a_i=(t_i,1-t_i)$ and $\theta_i=b_2+t_i(b_1-b_2)$.
Assign a size-$m$ subset $U$ to the first basis, put $s=\sum_{i\in U}t_i$
and $T=\sum_i t_i$. Its coefficient score is $m-T+2s$, maximized by
the top-$m$ sum. Removing the common translation and positive scale from
the weight SSE, minimizing it is equivalent to maximizing
$s^2+(T-s)^2$, hence $|s-T/2|$. The smallest and largest feasible sums
are the bottom-$m$ and top-$m$ sums; they are complementary and thus
correspond to the same unlabeled partitions. A convex quadratic attains
its maximum over these sums only at those endpoints, unless all feasible
sums coincide, in which case all partitions are tied. All endpoint
realizations are precisely the tied sorting optima, proving the claim.
\end{proof}

The condition on $b_1,b_2$ matters: identical bases make all weight
partitions optimal. Nor does the proposition certify a heuristic clustering
algorithm's output. For $m=1$ every unlabeled partition is the singleton
partition, independently of $K$.

\begin{proposition}[Isotropy on the simplex tangent space]
\label{prop:simplex-isotropy}
If $P_0BB^\top P_0=\lambda P_0$, $\lambda>0$, then
$J_\Theta(H)=\lambda J_A(H)$ for every balanced $H$.
\end{proposition}
\begin{proof}
For $X=(I-Q_H)A$, stochasticity gives $X\ones=0$, so $XP_0=X$.
Thus $\operatorname{tr}(XBB^\top X^\top)=\lambda\|X\|_F^2$.
\end{proof}

\subsection{Exact interior examples and sufficient information}
\label{app:decision-examples}

Set $A=N/23$, $B=[I_3\ \ 0_{3\times6}]$, where
\begin{equation}
 N=\begin{pmatrix}
 8&11&4\\8&5&10\\10&9&4\\5&11&7\\6&10&7\\11&5&7
 \end{pmatrix}. \label{eq:decision-counter-matrix}
\end{equation}
Every router entry is at least $4/23$, and the determinant of the first
three rows of $N$ is $276\ne0$; $B$ has three unit singular values.
The unique coefficient optimum has labels $(2,3,1,2,3,1)$, partition
$\mathcal C=\{\{1,4\},\{2,5\},\{3,6\}\}$, and score $60/23$.
For an explicit optimality certificate on the integer scale, subtract
column prices $(1,7/2,0)$ from each row of $N$. Each row now has a
\emph{strict} unique maximum at the assigned column, and those maxima
use each column twice. Adding the fixed price contribution to their sum
gives 60. Swapping the assignments of rows 4 and 5 gives score 59;
all scores are integers, so the runner-up gap is exactly $1/23$.

The unique weight and router-clustering optimum is instead
$\mathcal D=\{\{1,3\},\{2,6\},\{4,5\}\}$, cost $14/529$;
its runner-up costs $31/529$, while $\mathcal C$ costs $41/529$.
Table~\ref{tab:decision-certificates} exhausts the 15 unlabeled partitions.
Thus isotropic bases do not equate coefficient assignment with free-center
clustering. Strict gaps imply an open neighborhood of disagreement.
Replacing $A$ by $(1-r)\ones\ones^\top/3+rA$, $0<r\le1$, multiplies
coefficient-score differences by $r$ and clustering-cost differences by
$r^2$. The same example can therefore approach the uniform router;
its rank margin then shrinks. Proximity to uniformity alone is not an
agreement condition.

For $K=4$, append a fourth column and two rows: let $\widetilde N$ have
first six rows $(16N_{i1}+23,16N_{i2}+23,16N_{i3}+23,23)$ and last two
rows $(23,23,23,391)$. Then $\widetilde A=\widetilde N/460$ and
$\widetilde B=[I_4\ \ 0_{4\times8}]$ are positive and full rank.
Enumerating its 105 pair partitions gives $\mathcal C$ and $\mathcal D$
respectively, with the pair $\{7,8\}$ added. The $K=3$ example already
establishes non-equivalence; this extension matches the experimental
number of bases.

\paragraph{Weight--response tradeoff.}
In Equation~\eqref{eq:decision-counter-matrix}, replace $B$ by
$[\operatorname{Diag}(4,1,1)\ \ 0]$. Router clustering still uniquely
selects $\mathcal D$. Weight clustering uniquely selects
$\mathcal E=\{\{1,2\},\{3,6\},\{4,5\}\}$, with
\[
 (J_\Theta(\mathcal D),J_A(\mathcal D))=(119,14)/529,
 \qquad
 (J_\Theta(\mathcal E),J_A(\mathcal E))=(65,50)/529.
\]
For $\eta_B=1$, Theorem~\ref{thm:response-partition} gives response-error
increase $2Dm(50-14)/529=1296/529$ when choosing $\mathcal E$.
Neither objective uniformly controls the other.

\begin{corollary}[Effective weights do not determine the response optimum]
\label{cor:decision-information}
There exist two strictly positive full-rank routers and full-row-rank
bases with the same product but different unique balanced
Hilbert--Schmidt-response-optimal partitions.
\end{corollary}
\begin{proof}
Use the isotropic example above and
\[
 M=\tfrac1{10}\begin{pmatrix}2&1&7\\4&4&2\\5&2&3\end{pmatrix},
 \qquad A'=AM,\quad B'=M^{-1}B.
\]
The matrix is positive and row stochastic, with determinant $-7/100$.
Both factors retain their ranks, $A'$ is strictly positive, and
$A'B'=AB$. Router-clustering costs now have denominator $230^2$:
$\mathcal E$ uniquely costs $200/52900$, and runner-up $\mathcal D$
costs $212/52900$ (Table~\ref{tab:decision-certificates}). Theorem~\ref{thm:response-partition}
therefore selects $\mathcal D$ for the first factorization and
$\mathcal E$ for the second. Any rule observing only their common product
has the same input and cannot always return the unique optimum in both.
\end{proof}

Both worlds use their stated Euclidean factor metric. This is not a
passive coordinate substitution accompanied by a transformed optimizer.
The result does not make full reconstruction necessary: the
partition-dependent term in Theorem~\ref{thm:response-partition} is a
sum of within-group off-diagonal entries of $AA^\top$. Those entries
are already specified by off-diagonal response blocks. This observation
is consistent with the insufficiency of Gram information for recovering
all factors in Section~\ref{sec:response_identifiability}: deciding and
identifying are distinct targets.

\begin{table}[t]
\centering\small
\begin{tabular}{lrrrr}
\toprule
Partition & $N$ & $N\operatorname{Diag}(4,1,1)$ & $NM_{\rm int}$ & $N_0$ \\
\midrule
$12\,|\,34\,|\,56$ & 80 & 455 & 950 & 1070 \\
$12\,|\,35\,|\,46$ & 85 & 475 & 1015 & 636 \\
$12\,|\,36\,|\,45$ & 50 & 65 & 200 & 1358 \\
$13\,|\,24\,|\,56$ & 56 & 341 & 830 & 758 \\
$13\,|\,25\,|\,46$ & 59 & 389 & 917 & 732 \\
$13\,|\,26\,|\,45$ & 14 & 119 & 212 & 1190 \\
$14\,|\,23\,|\,56$ & 62 & 347 & 668 & 842 \\
$14\,|\,25\,|\,36$ & 41 & 146 & 347 & 842 \\
$14\,|\,26\,|\,35$ & 31 & 286 & 457 & 578 \\
$15\,|\,23\,|\,46$ & 71 & 401 & 809 & 636 \\
$15\,|\,24\,|\,36$ & 47 & 152 & 401 & 662 \\
$15\,|\,26\,|\,34$ & 35 & 320 & 533 & 806 \\
$16\,|\,23\,|\,45$ & 56 & 161 & 374 & 1202 \\
$16\,|\,24\,|\,35$ & 67 & 322 & 781 & 506 \\
$16\,|\,25\,|\,34$ & 65 & 350 & 803 & 914 \\
\bottomrule\end{tabular}
\caption{Complete integer SSE certificates. A label such as $12|34|56$ denotes three pairs. Divide columns by $23^2$, $23^2$, $230^2$, and $92^2$, respectively, for the stated router/effective-weight costs. Here $M_{\rm int}=10M$. Each column is computed by Equation~\eqref{eq:pairwise-sse}; no numerical optimizer is used.}
\label{tab:decision-certificates}
\end{table}


\subsection{A fixed loss that crosses a decision boundary}
\label{app:decision-crossing-proof}

For Theorem~\ref{thm:decision-crossing}, take $B_0=[I_3\ \ 0_{3\times6}]$
and $A_0=N_0/92$, with
\begin{equation}
 N_0=\begin{pmatrix}
 40&24&28\\24&44&24\\32&28&32\\43&37&12\\24&20&48\\37&31&24
 \end{pmatrix}. \label{eq:crossing-counter-matrix}
\end{equation}
Its rows sum to 92 and are positive. The first three rows are linearly
independent, so $A_0$ has column rank three. Define assignments
$p=(1,2,3,2,3,1)$ and $q=(1,2,3,1,3,2)$. They alone maximize the
coefficient score, both attaining $238/92$. Indeed, subtract integer
column prices $(6,0,0)$ from $N_0$. Rows 1,2,3,5 have strict maxima at
columns 1,2,3,3; only rows 4 and 6 tie between columns 1 and 2. Capacity
two forces one to each column, leaving exactly $p,q$. The next distinct
integer assignment score is 228; only positivity of the remaining gap
is needed. Table~\ref{tab:decision-certificates} shows that global weight
clustering uniquely selects the unlabeled partition of $p$, costing
$506/8464$ versus $578/8464$ for the runner-up.

Set $E=H_p-H_q$ with rows $e_i^\top$, and define fixed matrices by
\begin{equation}
 (G_0)_i=B_0^\top C_i^2e_i,\qquad
 T=A_0B_0-G_0,\qquad
 \ell(\Theta)=\tfrac12\|\Theta-T\|_F^2, \label{eq:crossing-target}
\end{equation}
where rows of $G_0$ are written as columns in the first expression and
$C_i$ is evaluated at $A_0$. Initialize the reference solution at
$Z_0=\log A_0$, $B_0$, at time zero. The loss is fixed, smooth and
bounded below; its gradient there is $G_0$. The exact chain rule is
\[
 \dot z_i=-C_iB_0g_i,\qquad
 \dot a_i=-C_i^2B_0g_i,\qquad
 \dot B=-A_0^\top G_0.
\]
Consequently Equation~\eqref{eq:crossing-speed} holds. Its strict sign
follows because $\ker C_i=\operatorname{span}(\ones)$, every nonzero
$e_i$ is orthogonal to $\ones$, and $B_0^\top$ is injective. At least
one $e_i$ is nonzero.

The smooth vector field has a solution on a two-sided neighborhood of
zero. By continuity, all other assignment gaps, the unique clustering
gap, positive entries and the two rank margins persist in a smaller
neighborhood. The strict negative derivative gives $S_p>S_q$ just
before zero and $S_p<S_q$ just after it. Weight clustering keeps $p$.
Choose a sufficiently small negative-time point as a forward initial
condition: the same fixed squared loss now trains from agreement into
disagreement. Choose two such endpoints with strict gaps. Continuous
dependence on initial conditions preserves their decisions and preserves
the weight optimum on the intervening compact interval for a small open
neighborhood of starting states. This proves Theorem~\ref{thm:decision-crossing}.

The numerical experiment uses an explicit negative tangent displacement
from $(Z_0,B_0)$ rather than inverse integration. Its observed crossing
checks that particular discrete initialization; the theorem's proof does
not rely on the numerical solver.

\paragraph{Local router velocities are not structurally restricted.}
At any positive router and full-row-rank basis, every row velocity $v_i$
with $\ones^\top v_i=0$ is realized at that point by choosing
\[
 g_i=-\frac{\tau^2}{\eta_Z}
 B^\top(BB^\top)^{-1}(C_i^2)^+v_i.
\]
Substitution into $\dot a_i=-(\eta_Z/\tau^2)C_i^2Bg_i$ gives $v_i$,
because $C_i^2(C_i^2)^+$ projects onto $\ones^\perp$. A fixed squared
loss with target $AB-G$ realizes that gradient at the point. The induced
$\dot B$ is not independently controlled. Natural tasks do not provide
arbitrary gradients, so this lemma supplies no typicality claim.

\subsection{Separation, margins, and trajectory error}
\label{app:decision-stability}

\begin{proposition}[A sufficient agreement regime]
\label{prop:near-hard-agreement}
Suppose a balanced labeling $c$ satisfies $a_{i,c(i)}\ge1-\epsilon_A$,
$0\le\epsilon_A<1/2$. Let $d$ be the minimum separation of distinct
basis rows and $D_B$ their maximum separation. If
$d>4\epsilon_A D_B$, coefficient assignment and global balanced weight
clustering uniquely select $c$.
\end{proposition}
\begin{proof}
Each own coefficient strictly exceeds the others and its pointwise best
labeling is balanced, so coefficient assignment is unique. Every
$\theta_i$ is within $\epsilon=\epsilon_A D_B$ of $b_{c(i)}$.
Within-group distances are at most $2\epsilon$; between-group distances
are at least $d-2\epsilon>2\epsilon$. All equal-capacity partitions
have the same number of within-group pairs. A different partition loses
and gains the same positive number of pairs; every new pair is strictly
more expensive than every lost pair. Equation~\eqref{eq:pairwise-sse}
therefore proves strict SSE optimality.
\end{proof}

More generally, for any given balanced partition, minimum between-group
squared distance exceeding maximum within-group squared distance certifies
its global SSE optimality by the same argument. Failure of this sufficient
condition is not evidence of non-optimality. Under the proposition's
separation, farthest-first seeding selects one center from each true group;
pointwise nearest-center assignment is then unique and balanced. Recomputed
centers stay within their respective $\epsilon$ balls, so balanced Lloyd
iterations retain that partition. This condition was not shown to hold
for most trained checkpoints.

\begin{proposition}[Finite-time margin protection]
\label{prop:decision-margin}
Initially let the unique labeled coefficient winner have gap $\Gamma_A>0$
and the unique unlabeled global weight-clustering winner have SSE gap
$\Gamma_\Theta>0$, with identical partitions. Let
$X_0=(I-\ones\ones^\top/L)\Theta(0)$. Both decisions remain unchanged
under perturbations satisfying
\begin{equation}
 \sqrt{2L}\|\Delta A\|_F<\Gamma_A,\qquad
 2\|X_0\|_F\|\Delta\Theta\|_F+\|\Delta\Theta\|_F^2<\Gamma_\Theta.
 \label{eq:decision-margin}
\end{equation}
\end{proposition}
\begin{proof}
The difference of two assignment matrices has norm at most $\sqrt{2L}$;
Cauchy--Schwarz bounds every score-gap perturbation. The difference of
two orthogonal projections has operator norm at most one. Expanding each
competing quadratic SSE difference bounds its change by
$2\|X_0\|_F\varepsilon+\varepsilon^2$, where
$\varepsilon=\|\Delta\Theta\|_F$. Centering is valid because every
$Q_H$ preserves $\ones$ and centering does not increase perturbation norm.
\end{proof}

If $\|B(t)\|_2\le\beta$ and $\|\nabla\ell(\Theta(t))\|_F\le M$
on an interval, then
\begin{equation}
 \|\dot A\|_F\le v_A:=\frac{\eta_Z\beta M}{4\tau^2},\qquad
 \|\dot\Theta\|_F\le v_\Theta:=
 \left(\eta_BL+\frac{\eta_Z\beta^2}{4\tau^2}\right)M.
 \label{eq:decision-speed-bounds}
\end{equation}
Indeed, Gershgorin's bound gives
$\|C_i\|_2\le\max_k2a_{ik}(1-a_{ik})\le1/2$, and
$\|A\|_2^2\le\|A\|_F^2\le L$.
Substituting $v_At$ and $v_\Theta t$ for the perturbations in
Equation~\eqref{eq:decision-margin} yields an explicit protected interval.
This is conditional stability, not attraction toward agreement.

\paragraph{Proof of Equation~\eqref{eq:trajectory-sharing-bound}.}
Write $E(t)=\Theta_s(t)-\Theta_h(t)$ and subtract the two ODEs:
\[
 \dot E=-(\mathcal R_s-\mathcal R_H)\nabla\ell(\Theta_s)
         -\mathcal R_H\bigl(\nabla\ell(\Theta_s)-\nabla\ell(\Theta_h)\bigr).
\]
The hypotheses give an upper Dini derivative of $\|E\|_F$ bounded by
$M\delta+c\|E\|_F$, since $\|HH^\top\|_2=m$ and
$c=\eta_BmL_\ell$. Integrating this scalar inequality yields
$\|E(t)\|_F\le e^{ct}\|E(0)\|_F+M\delta(e^{ct}-1)/c$;
the continuous $c=0$ limit is $\|E(0)\|_F+M\delta t$.
Group-mean initialization gives $\|E(0)\|_F=\sqrt{J_\Theta(H)}$.
The bound requires response control over the entire interval, not just
at its start, and supplies no lower bound or ranking of final task losses.
